
% 编译建议：命令行编译的话 XeLaTeX 连续编译两次，以正确生成目录、引用和总页数。 vscode 的 LaTeX Workshop 插件可直接使用“Build with recipe”。
\documentclass[UTF8,zihao=-4,a4paper,fontset=windows]{ctexart}

%==================== 基本信息 ====================
\newcommand{\StudentName}{Rpc}
\newcommand{\Major}{计算机科学与技术}
\newcommand{\ReportTitle}{机器学习入门知识总结与基本模型解析}
\newcommand{\ReportDate}{2026年7月29日}
\newcommand{\KnowledgeBaseCommit}{99ee977}

%==================== 宏包 ====================
\usepackage[a4paper,top=2.45cm,bottom=2.35cm,left=2.55cm,right=2.55cm,headheight=15pt]{geometry}
\usepackage[table]{xcolor}
\usepackage{amsmath,amssymb,mathtools,bm}
\usepackage{booktabs,longtable,tabularx,array,multirow}
\usepackage{enumitem}
\usepackage{float}
\usepackage{graphicx}
\usepackage[most]{tcolorbox}
\usepackage{tikz}
\usetikzlibrary{arrows.meta,positioning,shapes.geometric,calc}
\usepackage{fancyhdr}
\usepackage{lastpage}
\usepackage{hyperref}
\usepackage{microtype}
\usepackage{multicol}

%==================== 色彩与版式 ====================
\definecolor{Navy}{HTML}{16324F}
\definecolor{AcademicBlue}{HTML}{245A7A}
\definecolor{LightBlue}{HTML}{EDF5F8}
\definecolor{Accent}{HTML}{B56B28}
\definecolor{SoftGray}{HTML}{F4F6F7}
\definecolor{TextGray}{HTML}{4A5560}
\definecolor{FormulaGreen}{HTML}{2F6B55}

\hypersetup{
  colorlinks=true,
  linkcolor=AcademicBlue,
  citecolor=FormulaGreen,
  urlcolor=AcademicBlue,
  pdftitle={\ReportTitle},
  pdfauthor={\StudentName}
}
\Urlmuskip=0mu plus 1mu
\renewcommand{\contentsname}{目录}
\setlength{\parindent}{2em}
\setlength{\parskip}{0.22em}
\renewcommand{\baselinestretch}{1.30}
\setlist{itemsep=0.18em,topsep=0.35em,leftmargin=2.2em}
\setcounter{tocdepth}{2}
\setcounter{secnumdepth}{3}

\ctexset{
  part={
    format=\centering\Huge\bfseries\color{Navy},
    name={第,部分},
    number=\chinese{part},
    beforeskip=0.8em,
    afterskip=1.6em
  },
  section={
    format=\Large\bfseries\color{Navy},
    number=\chinese{section}、,
    beforeskip=1.15em,
    afterskip=0.65em
  },
  subsection={
    format=\large\bfseries\color{AcademicBlue},
    number=\arabic{section}.\arabic{subsection},
    beforeskip=0.9em,
    afterskip=0.45em
  },
  subsubsection={
    format=\normalsize\bfseries\color{TextGray},
    number=\arabic{section}.\arabic{subsection}.\arabic{subsubsection}
  }
}

% 目录采用三栏排版，缩小“部分”条目以避免在窄栏中断行。
\makeatletter
\renewcommand*\l@part[2]{%
  \ifnum\c@tocdepth>-2
    \addpenalty{-\@highpenalty}
    \vskip 0.35em
    \begingroup
      \parindent\z@
      \rightskip\@pnumwidth
      \parfillskip-\@pnumwidth
      \leavevmode\bfseries\color{AcademicBlue}
      #1\nobreak\hfill\nobreak
      \hb@xt@\@pnumwidth{\hss #2}\par
    \endgroup
  \fi
}
\makeatother

\pagestyle{fancy}
\fancyhf{}
\fancyhead[L]{\small\color{TextGray}机器学习理论与基本模型}
\fancyhead[R]{\small\color{TextGray}学术总结汇报}
\fancyfoot[C]{\small 第\thepage{}页，共\pageref{LastPage}页}
\renewcommand{\headrulewidth}{0.4pt}
\renewcommand{\footrulewidth}{0pt}

\newcolumntype{Y}{>{\raggedright\arraybackslash}X}
\newcolumntype{C}[1]{>{\centering\arraybackslash}p{#1}}
\newcolumntype{L}[1]{>{\raggedright\arraybackslash}p{#1}}

\newtcolorbox{definitionbox}[1]{
  enhanced,breakable,colback=LightBlue,colframe=AcademicBlue,
  boxrule=0.7pt,arc=1.5mm,left=2mm,right=2mm,top=1.4mm,bottom=1.4mm,
  title=#1,fonttitle=\bfseries,coltitle=Navy
}
\newtcolorbox{derivationbox}[1]{
  enhanced,breakable,colback=white,colframe=FormulaGreen,
  borderline west={2.5pt}{0pt}{FormulaGreen},boxrule=0.35pt,arc=1mm,
  left=2mm,right=2mm,top=1.3mm,bottom=1.3mm,
  title=#1,fonttitle=\bfseries,coltitle=FormulaGreen
}
\newtcolorbox{notebox}[1]{
  enhanced,breakable,colback=SoftGray,colframe=Accent,
  borderline west={2.5pt}{0pt}{Accent},boxrule=0.35pt,arc=1mm,
  left=2mm,right=2mm,top=1.3mm,bottom=1.3mm,
  title=#1,fonttitle=\bfseries,coltitle=Accent
}

\DeclareMathOperator*{\argmin}{arg\,min}
\DeclareMathOperator*{\argmax}{arg\,max}
\DeclareMathOperator{\diag}{diag}
\DeclareMathOperator{\sign}{sign}
\DeclareMathOperator{\tr}{tr}
\newcommand{\R}{\mathbb{R}}
\newcommand{\E}{\mathbb{E}}
\newcommand{\Prob}{\mathbb{P}}
\newcommand{\norm}[1]{\left\lVert #1\right\rVert}
\newcommand{\abs}[1]{\left|#1\right|}
\newcommand{\dataset}{\mathcal{D}}
\newcommand{\loss}{\mathcal{L}}

\begin{document}

%==================== 封面 ====================
\begin{titlepage}
  \thispagestyle{empty}
  \begin{tikzpicture}[remember picture,overlay]
    \fill[Navy] (current page.north west) rectangle ([yshift=-3.4cm]current page.north east);
    \fill[AcademicBlue!13] ([yshift=2.0cm]current page.south west) rectangle (current page.south east);
    \draw[AcademicBlue,line width=1.1pt]
      ([xshift=2.5cm,yshift=-4.0cm]current page.north west)--
      ([xshift=-2.5cm,yshift=-4.0cm]current page.north east);
  \end{tikzpicture}

  \vspace*{0.95cm}
  {\color{white}\Large\bfseries 暑期学习作业·第二部分}\par
  \vspace{3.1cm}
  \begin{center}
    {\Huge\bfseries\color{Navy}\ReportTitle\par}
    \vspace{0.55cm}
    {\LARGE\color{AcademicBlue}---从学习范式、优化与评估到经典模型的数学基础\par}
    \vspace{1.6cm}

    \begin{tcolorbox}[
      width=0.86\textwidth,colback=LightBlue,colframe=AcademicBlue,
      arc=2mm,boxrule=0.8pt,left=4mm,right=4mm,top=3mm,bottom=3mm
    ]
      \centering\large
      第一部分总结机器学习、深度学习与强化学习的统一学习框架，\par
      第二部分按照知识库提纲推导并比较经典模型。
    \end{tcolorbox}

    \vfill
    \renewcommand{\arraystretch}{1.65}
    \begin{tabular}{>{\bfseries\color{Navy}}r l}
      姓名：& \StudentName \\
      专业：& \Major \\
      资料来源：& AI-Knowledge-Base（提交 \texttt{\KnowledgeBaseCommit}）\\
      完成日期：& \ReportDate
    \end{tabular}
    \vfill
  \end{center}
\end{titlepage}

%==================== 摘要与目录 ====================
\pagenumbering{Roman}
\begin{definitionbox}{摘要}
本文从理论与数学层面对人工智能入门知识进行系统总结。第一部分以经验风险最小化为统一视角，
区分机器学习、深度学习和强化学习三种学习范式，说明数据划分、前向传播、损失计算、
反向传播、参数优化与收敛判断之间的关系，并汇总回归、分类、聚类和强化学习中的基础评价指标。
第二部分依据 AI-Knowledge-Base 的学习路径，从线性回归开始，依次解析逻辑回归、
决策树、朴素贝叶斯、支持向量机、随机森林、梯度提升、K-Means、层次聚类、主成分分析、
多层感知机、卷积神经网络和自注意力机制。对每类模型均给出问题定义、核心公式、
优化目标、适用条件与主要局限。本文的目的不是罗列算法名称，而是建立
“任务定义--模型假设--目标函数--求解方法--评价指标--泛化能力”的统一分析框架，
为后续代码复现和研究方向学习提供理论基础。

\textbf{关键词：}机器学习；深度学习；强化学习；梯度优化；模型评估；经典算法
\end{definitionbox}

\vspace{0.45em}
\begin{multicols}{3}
  \scriptsize
  \tableofcontents
\end{multicols}
\clearpage
\pagenumbering{arabic}

%====================================================
\part{基本原理}
%====================================================

\noindent
本文的模型范围和学习顺序以 AI-Knowledge-Base 的现有目录为主要提纲，
并针对作业要求补充统一训练框架、收敛条件、评价指标和必要的数学推导\cite{aikb}。

\section{统一的问题表述与数学基础}

\subsection{从数据到学习问题}

机器学习的基本目标，是利用有限样本学习一个可以推广到未见数据的映射。
设训练集为
\begin{equation}
  \dataset=\{(\bm{x}_i,y_i)\}_{i=1}^{n},\qquad
  \bm{x}_i\in\R^d ,
\end{equation}
其中 $\bm{x}_i$ 表示第 $i$ 个样本的 $d$ 维特征，$y_i$ 为监督信号。
模型记为 $f_{\bm{\theta}}$，参数为 $\bm{\theta}$。给定单样本损失
$\ell(f_{\bm{\theta}}(\bm{x}_i),y_i)$，训练通常写成正则化经验风险最小化问题：
\begin{equation}
  \hat{\bm{\theta}}
  =\argmin_{\bm{\theta}}
  \underbrace{\frac{1}{n}\sum_{i=1}^{n}
  \ell\!\left(f_{\bm{\theta}}(\bm{x}_i),y_i\right)}_{\text{经验风险}}
  +\lambda\underbrace{\Omega(\bm{\theta})}_{\text{复杂度约束}} .
  \label{eq:erm}
\end{equation}
式 \eqref{eq:erm} 把多数监督学习模型放在同一个框架内：
模型结构决定 $f_{\bm{\theta}}$ 的表达能力，损失函数定义“错误”，
优化器寻找较优参数，正则项控制复杂度。机器学习的核心并非只让训练误差最小，
而是使总体风险
\begin{equation}
  R(\bm{\theta})
  =\E_{(\bm{x},y)\sim p_{\mathrm{data}}}
  \left[\ell(f_{\bm{\theta}}(\bm{x}),y)\right]
\end{equation}
尽可能小。由于真实数据分布 $p_{\mathrm{data}}$ 未知，只能用有限样本近似，
所以产生了泛化问题\cite{mitchell}。

\begin{definitionbox}{训练集、验证集与测试集}
\begin{itemize}
  \item \textbf{训练集}用于拟合参数 $\bm{\theta}$；
  \item \textbf{验证集}用于选择超参数、模型结构和停止时刻；
  \item \textbf{测试集}只用于最终无偏评估，不能参与预处理参数估计或模型选择。
\end{itemize}
若在全部数据上先做标准化、特征选择或重复样本清理，再划分训练集与测试集，
测试信息可能泄漏到训练过程，最终指标会过于乐观。
\end{definitionbox}

\subsection{基础数学工具}

机器学习中的样本、参数和中间表示通常写成向量或矩阵。
对矩阵 $\bm{X}\in\R^{n\times d}$，第 $i$ 行为一个样本，
第 $j$ 列为一个特征。常用范数为
\begin{align}
  \norm{\bm{x}}_1 &= \sum_{j=1}^{d}\abs{x_j}, &
  \norm{\bm{x}}_2 &= \sqrt{\sum_{j=1}^{d}x_j^2}, &
  \norm{\bm{x}}_\infty &= \max_j\abs{x_j}.
\end{align}
$L_1$ 范数常用于产生稀疏参数，$L_2$ 范数常用于平滑地限制参数规模。
梯度
\begin{equation}
  \nabla_{\bm{\theta}}\loss
  =\left[
  \frac{\partial\loss}{\partial\theta_1},
  \ldots,
  \frac{\partial\loss}{\partial\theta_p}
  \right]^{\mathsf T}
\end{equation}
给出损失增长最快的局部方向，因此负梯度是最基本的下降方向。
概率论则用于描述不确定性：先验 $p(\theta)$、似然 $p(\dataset\mid\theta)$
和后验 $p(\theta\mid\dataset)$ 由贝叶斯公式联系：
\begin{equation}
  p(\theta\mid\dataset)
  =\frac{p(\dataset\mid\theta)p(\theta)}{p(\dataset)}.
\end{equation}

\section{机器学习、深度学习与强化学习}

\subsection{机器学习}

机器学习关注如何从数据中自动获得预测或决策规则\cite{mitchell}。
按照监督信号的形式，可划分为以下几类：

\begin{table}[H]
\centering
\caption{主要学习范式比较}
\begin{tabularx}{\textwidth}{L{2.6cm}Y Y Y}
\toprule
\textbf{范式} & \textbf{数据形式} & \textbf{目标} & \textbf{典型方法}\\
\midrule
监督学习 & $(\bm{x},y)$ 成对出现 & 学习输入到标签的映射 & 回归、分类、排序\\
无监督学习 & 只有 $\bm{x}$ & 发现结构、密度或低维表示 & 聚类、PCA、密度估计\\
半监督学习 & 少量有标签与大量无标签数据 & 同时利用两类数据 & 伪标签、一致性正则化\\
自监督学习 & 从数据自身构造监督信号 & 学习可迁移表示 & 掩码预测、对比学习\\
\bottomrule
\end{tabularx}
\end{table}

监督学习又分为回归与分类。回归输出连续值，分类输出离散类别或类别概率。
无监督学习没有唯一的“正确标签”，因此评价必须结合内部结构、外部标签或下游任务。

\subsection{深度学习}

深度学习是机器学习的一类方法，其核心是使用多层可微模块逐层学习表示\cite{goodfellow}。
第 $\ell$ 层可写为
\begin{equation}
  \bm{z}^{(\ell)}
  =\bm{W}^{(\ell)}\bm{a}^{(\ell-1)}+\bm{b}^{(\ell)},
  \qquad
  \bm{a}^{(\ell)}=\phi^{(\ell)}\!\left(\bm{z}^{(\ell)}\right).
  \label{eq:layer}
\end{equation}
与依赖人工特征的传统流程相比，深度网络可以把特征提取和任务预测放在同一目标函数中联合优化。
“深”并不只是层数增加，还意味着模型通过组合简单非线性变换形成层次化表示。

\begin{notebox}{机器学习与深度学习的关系}
深度学习不是与机器学习并列的完全独立领域，而是机器学习的一个重要分支。
传统模型通常参数较少、可解释性较强、适合中小规模结构化数据；
深度模型表示能力强，适合图像、语音和文本等高维数据，但往往需要更多数据、算力和正则化。
\end{notebox}

\subsection{强化学习}

强化学习研究智能体如何通过与环境交互来最大化长期回报\cite{sutton}。
马尔可夫决策过程表示为五元组
\begin{equation}
  \mathcal{M}=(\mathcal{S},\mathcal{A},P,R,\gamma),
\end{equation}
其中 $\mathcal{S}$ 为状态空间，$\mathcal{A}$ 为动作空间，
$P(s'\mid s,a)$ 为转移概率，$R(s,a)$ 为即时奖励，
$\gamma\in[0,1)$ 为折扣因子。轨迹从时刻 $t$ 开始的回报为
\begin{equation}
  G_t=\sum_{k=0}^{\infty}\gamma^k r_{t+k+1}.
\end{equation}
策略 $\pi(a\mid s)$ 描述在状态 $s$ 下选择动作 $a$ 的概率。
动作价值函数满足 Bellman 期望方程
\begin{equation}
  Q^\pi(s,a)
  =\E_\pi\!\left[
  r_{t+1}+\gamma Q^\pi(s_{t+1},a_{t+1})
  \mid s_t=s,a_t=a
  \right].
\end{equation}
经典 Q-Learning 直接逼近最优动作价值：
\begin{equation}
  Q(s_t,a_t)\leftarrow Q(s_t,a_t)
  +\alpha\left[
  r_{t+1}+\gamma\max_{a'}Q(s_{t+1},a')
  -Q(s_t,a_t)
  \right].
\end{equation}
方括号中的量称为时序差分误差。与监督学习不同，强化学习的监督信号由交互产生，
样本通常非独立同分布，还必须处理探索与利用的矛盾。

\section{模型训练、传播与优化}

\subsection{完整训练过程}

模型训练需要把数据处理、目标函数、数值优化和泛化评估放在同一流程中理解
\cite{bishop,hastie}。

\begin{center}
\begin{tikzpicture}[
  node distance=7mm and 5mm,
  >=Latex,
  box/.style={rounded corners=1.5mm,draw=AcademicBlue,fill=LightBlue,
              minimum width=3.0cm,minimum height=1.0cm,align=center},
  arrow/.style={->,thick,AcademicBlue}
]
  \node[box] (data) {数据准备\\划分与预处理};
  \node[box,right=of data] (forward) {前向传播\\得到预测};
  \node[box,right=of forward] (loss) {计算损失\\与评价指标};
  \node[box,below=of loss] (backward) {反向传播\\计算梯度};
  \node[box,left=of backward] (update) {优化器更新\\模型参数};
  \node[box,left=of update] (validate) {验证与诊断\\收敛或早停};
  \draw[arrow] (data)--(forward);
  \draw[arrow] (forward)--(loss);
  \draw[arrow] (loss)--(backward);
  \draw[arrow] (backward)--(update);
  \draw[arrow] (update)--(validate);
  \draw[arrow] (validate.west) -- ++(-0.55,0) |- (forward.south);
\end{tikzpicture}
\end{center}

一次完整训练包括：
\begin{enumerate}
  \item 固定数据版本并划分训练集、验证集和测试集；
  \item 仅在训练集上拟合标准化或特征处理参数；
  \item 初始化模型参数，将一个批次输入模型完成前向传播；
  \item 根据预测和真实目标计算损失；
  \item 通过反向传播求梯度，由优化器更新参数；
  \item 重复多个批次与轮次，并监控训练集和验证集曲线；
  \item 根据预先约定的准则选择模型，最后在测试集上评估一次。
\end{enumerate}
遍历一次训练集称为一个 \textit{epoch}；一次用于计算梯度的样本集合称为
\textit{mini-batch}；一次参数更新称为一个 \textit{iteration}。

\subsection{什么是前向传播}

前向传播是在参数固定时，把输入按计算图方向逐层变换成预测值的过程。
对于 $L$ 层网络，令 $\bm{a}^{(0)}=\bm{x}$，则
\begin{align}
  \bm{z}^{(\ell)}
  &=\bm{W}^{(\ell)}\bm{a}^{(\ell-1)}+\bm{b}^{(\ell)},\\
  \bm{a}^{(\ell)}
  &=\phi^{(\ell)}(\bm{z}^{(\ell)}),
  \qquad \ell=1,\ldots,L.
\end{align}
最终输出 $\hat{\bm{y}}=\bm{a}^{(L)}$，损失为
$\loss=\ell(\hat{\bm{y}},\bm{y})$。前向传播只负责计算预测和损失，
本身不改变参数；模型真正“学习”发生在梯度计算与参数更新阶段。

\subsection{反向传播与链式法则}

反向传播不是一种独立的优化器，而是高效计算复合函数梯度的方法\cite{rumelhart}。
定义第 $\ell$ 层的误差信号
\begin{equation}
  \bm{\delta}^{(\ell)}
  =\frac{\partial\loss}{\partial\bm{z}^{(\ell)}}.
\end{equation}
由链式法则可递推得到
\begin{align}
  \bm{\delta}^{(L)}
  &=\nabla_{\bm{a}^{(L)}}\loss
    \odot \phi^{(L)\prime}(\bm{z}^{(L)}),\\
  \bm{\delta}^{(\ell)}
  &=\left(\bm{W}^{(\ell+1)}\right)^{\mathsf T}
    \bm{\delta}^{(\ell+1)}
    \odot \phi^{(\ell)\prime}(\bm{z}^{(\ell)}),\\
  \frac{\partial\loss}{\partial\bm{W}^{(\ell)}}
  &=\bm{\delta}^{(\ell)}
    \left(\bm{a}^{(\ell-1)}\right)^{\mathsf T},\qquad
  \frac{\partial\loss}{\partial\bm{b}^{(\ell)}}
  =\bm{\delta}^{(\ell)}.
\end{align}
$\odot$ 表示逐元素乘法。这个递推把输出误差逐层传回，同时复用前向传播保存的中间量，
比对每个参数分别求导高效得多。

\subsection{梯度下降及其变体}

最基本的参数更新为
\begin{equation}
  \bm{\theta}_{t+1}
  =\bm{\theta}_{t}-\eta_t\nabla_{\bm{\theta}}\loss(\bm{\theta}_t),
  \label{eq:gd}
\end{equation}
其中 $\eta_t$ 为学习率。批量梯度下降使用全部训练样本，
随机梯度下降每次使用一个样本，小批量梯度下降在计算效率和梯度方差之间折中。
带动量的更新为
\begin{align}
  \bm{v}_{t+1}&=\beta\bm{v}_t+\nabla\loss(\bm{\theta}_t),\\
  \bm{\theta}_{t+1}&=\bm{\theta}_t-\eta\bm{v}_{t+1}.
\end{align}
Adam 同时估计梯度的一阶矩和二阶矩，并进行偏差修正\cite{kingma}：
\begin{align}
  \bm{m}_t&=\beta_1\bm{m}_{t-1}+(1-\beta_1)\bm{g}_t,\\
  \bm{v}_t&=\beta_2\bm{v}_{t-1}+(1-\beta_2)\bm{g}_t^2,\\
  \bm{\theta}_{t+1}
  &=\bm{\theta}_t-\eta
  \frac{\hat{\bm{m}}_t}{\sqrt{\hat{\bm{v}}_t}+\epsilon}.
\end{align}
学习率过大可能导致震荡或发散，过小则会使收敛缓慢。
因此需要结合预热、分段衰减、余弦退火或验证集早停等策略。

\subsection{什么是收敛}

\begin{definitionbox}{数学收敛与训练收敛}
若参数序列存在极限 $\bm{\theta}^{\star}$，且
$\lim_{t\to\infty}\norm{\bm{\theta}_t-\bm{\theta}^{\star}}=0$，
称参数序列收敛。实际训练无法无限迭代，通常使用以下近似判据：
\[
  \abs{\loss_{t+1}-\loss_t}<\varepsilon_L,\qquad
  \norm{\nabla\loss(\bm{\theta}_t)}_2<\varepsilon_g,\qquad
  \norm{\bm{\theta}_{t+1}-\bm{\theta}_t}_2<\varepsilon_\theta .
\]
也可以在验证指标连续若干轮不再改善时早停。
\end{definitionbox}

对于凸目标，满足适当平滑性和学习率条件时，梯度下降可以收敛到全局最优解；
对于深度网络的非凸目标，通常只能期望到达梯度较小的驻点或平坦区域。
小批量噪声还会使单步损失波动，因此不能仅凭一两次迭代判断是否收敛。

\begin{notebox}{收敛不等于泛化}
训练损失稳定下降只说明优化过程趋于稳定，不代表模型能处理新数据。
若训练损失继续下降而验证损失上升，模型正在过拟合。
因此收敛判断必须同时观察训练曲线、验证曲线、泛化间隔和多次独立运行的稳定性。
\end{notebox}

\subsection{欠拟合、过拟合与正则化}

欠拟合表示模型无法充分解释训练数据，训练误差和验证误差都较高；
过拟合表示训练误差很低但验证误差明显更高。常用控制方法包括：
\begin{itemize}
  \item \textbf{$L_2$ 正则化：}
  $\Omega(\bm{\theta})=\frac{1}{2}\norm{\bm{\theta}}_2^2$，
  使参数平滑缩小；
  \item \textbf{$L_1$ 正则化：}
  $\Omega(\bm{\theta})=\norm{\bm{\theta}}_1$，
  倾向产生稀疏参数；
  \item \textbf{早停：}在验证性能不再改善时停止训练；
  \item \textbf{数据增强：}用保持标签语义的变换扩大训练分布；
  \item \textbf{Dropout：}训练时随机屏蔽部分神经元，减少共适应；
  \item \textbf{交叉验证：}在有限样本下更稳定地估计泛化性能。
\end{itemize}

\section{基础衡量指标}

\subsection{回归指标}

设真实值为 $y_i$，预测值为 $\hat y_i$，样本均值为 $\bar y$。
常用指标如下：
\begin{align}
  \mathrm{MAE}
  &=\frac{1}{n}\sum_{i=1}^{n}\abs{y_i-\hat y_i},\\
  \mathrm{MSE}
  &=\frac{1}{n}\sum_{i=1}^{n}(y_i-\hat y_i)^2,\\
  \mathrm{RMSE}
  &=\sqrt{\mathrm{MSE}},\\
  R^2
  &=1-\frac{\sum_i(y_i-\hat y_i)^2}
  {\sum_i(y_i-\bar y)^2}.
\end{align}
MAE 对异常值更稳健；MSE 对大误差惩罚更强；RMSE 与目标变量量纲一致；
$R^2$ 衡量相对于均值基线减少了多少平方误差。
在测试集上，若模型比均值基线更差，$R^2$ 可以小于 $0$，并非总在 $[0,1]$ 内。

\subsection{分类指标}

二分类混淆矩阵包含真正例 $TP$、假正例 $FP$、真负例 $TN$ 和假负例 $FN$。
\begin{table}[H]
\centering
\caption{二分类混淆矩阵}
\begin{tabular}{c cc}
\toprule
& \textbf{预测为正} & \textbf{预测为负}\\
\midrule
\textbf{实际为正} & $TP$ & $FN$\\
\textbf{实际为负} & $FP$ & $TN$\\
\bottomrule
\end{tabular}
\end{table}

\begin{align}
  \mathrm{Accuracy}
  &=\frac{TP+TN}{TP+TN+FP+FN},\\
  \mathrm{Precision}
  &=\frac{TP}{TP+FP},\\
  \mathrm{Recall}
  &=\frac{TP}{TP+FN},\\
  \mathrm{Specificity}
  &=\frac{TN}{TN+FP},\\
  F_1
  &=\frac{2\cdot\mathrm{Precision}\cdot\mathrm{Recall}}
  {\mathrm{Precision}+\mathrm{Recall}}.
\end{align}
Precision 回答“预测为正的样本中有多少是真的”，Recall 回答
“实际为正的样本中找回了多少”。类别不平衡时，Accuracy 可能掩盖少数类错误，
应同时报告 Precision、Recall、$F_1$、PR-AUC 或 ROC-AUC。
多分类任务还需说明采用 macro、micro 还是 weighted 平均。

对于概率预测，二分类交叉熵为
\begin{equation}
  \mathrm{BCE}
  =-\frac{1}{n}\sum_{i=1}^{n}
  \left[y_i\log p_i+(1-y_i)\log(1-p_i)\right].
\end{equation}
它评价概率分布而不只是最终标签，对“高置信度但错误”的预测惩罚很大。

\subsection{聚类与降维指标}

聚类没有统一标签时，可以使用轮廓系数：
\begin{equation}
  s(i)=\frac{b(i)-a(i)}{\max\{a(i),b(i)\}}\in[-1,1],
\end{equation}
其中 $a(i)$ 为样本到同簇样本的平均距离，$b(i)$ 为到最近其他簇的平均距离。
也可报告 Calinski--Harabasz 指数和 Davies--Bouldin 指数。
这些内部指标反映紧致性与分离性，但不一定对应具体任务价值。
PCA 通常用累计解释方差比与重建误差评价，后文将给出公式。

\subsection{强化学习指标}

强化学习最直接的指标是多个独立回合的平均累计回报：
\begin{equation}
  \bar G=\frac{1}{M}\sum_{m=1}^{M}G^{(m)}.
\end{equation}
还应报告成功率、回报方差、达到目标性能所需的环境交互次数（样本效率）、
训练时间以及不同随机种子下的稳定性。若只报告最优一次运行，
容易高估算法真实性能。

\section{实验设计与可信结论}

一个可信实验不仅需要正确实现算法，还要建立完整证据链\cite{bishop,hastie}：
\begin{enumerate}
  \item \textbf{明确问题与假设：}输入、输出、数据分布和评价目标必须清晰；
  \item \textbf{设置公平 Baseline：}保持数据、划分、指标、训练预算和调参条件一致；
  \item \textbf{进行消融实验：}一次改变一个关键因素，验证各模块的独立贡献；
  \item \textbf{检查稳定性：}使用多个随机种子，报告均值与标准差；
  \item \textbf{开展错误分析：}按类别、难度、数据来源和置信度分析失败样本；
  \item \textbf{记录可复现信息：}保存代码版本、环境、配置、日志和模型选择规则。
\end{enumerate}
统计显著不等于实际重要，指标提高也不自动证明理论解释正确。
结论应区分观察事实、机制解释和适用范围。

%====================================================
\part{基本模型解析}
%====================================================

\section{监督学习基础模型}

\subsection{线性回归}

\subsubsection{模型与最小二乘目标}

线性回归是知识库模型路径的起点\cite{aikb}。多元线性模型写成
\begin{equation}
  y_i=\beta_0+\bm{x}_i^{\mathsf T}\bm{\beta}+\varepsilon_i,
\end{equation}
把常数列并入设计矩阵后，有
\begin{equation}
  \bm{y}=\bm{X}\bm{\beta}+\bm{\varepsilon}.
\end{equation}
普通最小二乘法最小化残差平方和：
\begin{equation}
  J(\bm{\beta})
  =\frac{1}{2n}\norm{\bm{y}-\bm{X}\bm{\beta}}_2^2.
\end{equation}
对参数求梯度：
\begin{equation}
  \nabla_{\bm{\beta}}J
  =\frac{1}{n}\bm{X}^{\mathsf T}
  (\bm{X}\bm{\beta}-\bm{y}).
\end{equation}
令梯度为零得到正规方程
\begin{equation}
  \bm{X}^{\mathsf T}\bm{X}\hat{\bm{\beta}}
  =\bm{X}^{\mathsf T}\bm{y}.
\end{equation}
若 $\bm{X}^{\mathsf T}\bm{X}$ 可逆，则
\begin{equation}
  \hat{\bm{\beta}}
  =(\bm{X}^{\mathsf T}\bm{X})^{-1}
  \bm{X}^{\mathsf T}\bm{y}.
  \label{eq:ols}
\end{equation}
实际计算通常使用 QR 分解、SVD 或伪逆，而不是显式求逆，以提高数值稳定性。

\begin{derivationbox}{几何解释}
正规方程等价于
\[
  \bm{X}^{\mathsf T}(\bm{y}-\bm{X}\hat{\bm{\beta}})=\bm{0}.
\]
因此残差向量与 $\bm{X}$ 的列空间正交，
$\hat{\bm{y}}=\bm{X}\hat{\bm{\beta}}$ 是 $\bm{y}$ 在该列空间上的正交投影。
\end{derivationbox}

\subsubsection{假设、正则化与局限}

线性回归假设条件均值近似线性。误差独立、同方差等条件影响统计推断；
误差正态性不是求得 OLS 解或保证无偏性的必要条件，但常用于有限样本置信区间和显著性检验。
多重共线性会使参数方差变大。Ridge 回归和 Lasso 分别求解
\begin{align}
  \hat{\bm{\beta}}_{\mathrm{ridge}}
  &=\argmin_{\bm{\beta}}
  J(\bm{\beta})+\lambda\norm{\bm{\beta}}_2^2,\\
  \hat{\bm{\beta}}_{\mathrm{lasso}}
  &=\argmin_{\bm{\beta}}
  J(\bm{\beta})+\lambda\norm{\bm{\beta}}_1.
\end{align}
前者稳定参数，后者可以产生稀疏解。线性回归可解释、训练快，
但对异常值、非线性关系和强共线性敏感。

\subsection{逻辑回归}

\subsubsection{概率模型与对数几率}

逻辑回归用于二分类。它先构造线性打分
$z=\bm{w}^{\mathsf T}\bm{x}+b$，再用 Sigmoid 映射到概率：
\begin{equation}
  p(y=1\mid\bm{x})
  =\sigma(z)=\frac{1}{1+\exp(-z)}.
\end{equation}
对概率取对数几率得到
\begin{equation}
  \log\frac{p}{1-p}
  =\bm{w}^{\mathsf T}\bm{x}+b,
\end{equation}
所以逻辑回归假设的是“对数几率关于特征线性”，而非类别标签本身线性。
判定阈值不必固定为 $0.5$，应根据误报和漏报代价选择。

\subsubsection{最大似然与梯度}

在样本条件独立假设下，似然为
\begin{equation}
  L(\bm{w},b)
  =\prod_{i=1}^{n}
  p_i^{y_i}(1-p_i)^{1-y_i}.
\end{equation}
最大化似然等价于最小化负对数似然：
\begin{equation}
  J(\bm{w},b)
  =-\frac{1}{n}\sum_{i=1}^{n}
  \left[y_i\log p_i+(1-y_i)\log(1-p_i)\right].
\end{equation}
若把偏置并入参数，梯度可写成
\begin{equation}
  \nabla_{\bm{w}}J
  =\frac{1}{n}\bm{X}^{\mathsf T}
  (\bm{p}-\bm{y}).
\end{equation}
该目标对参数是凸函数，加入凸正则项后仍可用梯度法稳定求解。
多分类时使用 Softmax：
\begin{equation}
  p(y=k\mid\bm{x})
  =\frac{\exp(z_k)}{\sum_{j=1}^{K}\exp(z_j)}.
\end{equation}
逻辑回归输出可解释概率，适合作为强 Baseline；
但其决策边界在原始特征空间中是线性的，需要特征工程或核方法表达复杂边界。

\subsection{决策树}

\subsubsection{递归划分与纯度}

决策树通过递归选择特征和切分点，把输入空间划分为若干区域。
分类树常用熵、信息增益和基尼不纯度。若节点数据集 $D$ 中第 $k$ 类比例为 $p_k$，
则
\begin{align}
  H(D)&=-\sum_{k=1}^{K}p_k\log_2 p_k,\\
  \mathrm{Gini}(D)&=1-\sum_{k=1}^{K}p_k^2.
\end{align}
属性 $a$ 把 $D$ 划分为 $\{D^v\}_{v=1}^{V}$，信息增益为
\begin{equation}
  \mathrm{Gain}(D,a)
  =H(D)-\sum_{v=1}^{V}\frac{\abs{D^v}}{\abs D}H(D^v).
\end{equation}
信息增益偏向取值较多的属性，增益比用属性固有值进行归一化：
\begin{align}
  \mathrm{IV}(a)
  &=-\sum_{v=1}^{V}\frac{\abs{D^v}}{\abs D}
  \log_2\frac{\abs{D^v}}{\abs D},\\
  \mathrm{GainRatio}(D,a)
  &=\frac{\mathrm{Gain}(D,a)}{\mathrm{IV}(a)}.
\end{align}
CART 分类树通常选择加权基尼不纯度最小的切分。

\subsubsection{回归树与剪枝}

回归树在每个叶节点输出区域均值。对候选区域 $R_1,R_2$，切分目标为
\begin{equation}
  \min_{R_1,R_2}
  \left[
  \sum_{\bm{x}_i\in R_1}(y_i-\bar y_{R_1})^2+
  \sum_{\bm{x}_i\in R_2}(y_i-\bar y_{R_2})^2
  \right].
\end{equation}
树过深容易学习噪声。预剪枝通过最大深度、最小叶节点样本数或最小增益提前停止；
后剪枝先生成较大树，再依据验证误差或代价复杂度目标
\begin{equation}
  R_\alpha(T)=R(T)+\alpha\abs{T_{\mathrm{leaf}}}
\end{equation}
删除收益不足的分支。决策树能够处理非线性与特征交互，解释直观，
但对样本扰动敏感、方差较高，这也是随机森林等集成模型的出发点。

\subsection{朴素贝叶斯}

\subsubsection{贝叶斯决策与条件独立}

对类别 $C_k$ 和特征 $\bm{x}=(x_1,\ldots,x_d)$，贝叶斯公式给出
\begin{equation}
  p(C_k\mid\bm{x})
  =\frac{p(\bm{x}\mid C_k)p(C_k)}{p(\bm{x})}.
\end{equation}
朴素贝叶斯作出强条件独立假设：
\begin{equation}
  p(\bm{x}\mid C_k)
  =\prod_{j=1}^{d}p(x_j\mid C_k).
\end{equation}
由于分母对所有类别相同，最大后验决策为
\begin{equation}
  \hat y
  =\argmax_k\left[
  \log p(C_k)+
  \sum_{j=1}^{d}\log p(x_j\mid C_k)
  \right].
\end{equation}
使用对数可把连乘转换为求和并避免数值下溢。

\subsubsection{平滑与常见分布}

离散特征若某个计数为零，会使整个类别后验变为零。
对词频模型可使用加性平滑：
\begin{equation}
  \hat p(w_j\mid C_k)
  =\frac{N_{jk}+\alpha}
  {\sum_{r=1}^{V}N_{rk}+\alpha V},
\end{equation}
其中 $V$ 为词表大小，$\alpha=1$ 时为拉普拉斯平滑。
高斯朴素贝叶斯假设连续特征满足
\begin{equation}
  p(x_j\mid C_k)
  =\frac{1}{\sqrt{2\pi\sigma_{jk}^2}}
  \exp\left[-\frac{(x_j-\mu_{jk})^2}{2\sigma_{jk}^2}\right].
\end{equation}
多项式模型适合计数特征，伯努利模型适合二值出现特征。
朴素贝叶斯训练和预测都很快，在文本分类与小样本任务中常表现良好；
其主要局限是条件独立假设通常不完全成立，概率校准也可能较差。

\section{进阶监督学习与集成方法}

\subsection{支持向量机}

\subsubsection{最大间隔与硬间隔}

对 $y_i\in\{-1,+1\}$，线性分类超平面为
$\bm{w}^{\mathsf T}\bm{x}+b=0$。规范化约束后，硬间隔 SVM 求解
\begin{equation}
  \min_{\bm{w},b}\frac{1}{2}\norm{\bm{w}}_2^2
  \quad
  \text{s.t.}\quad
  y_i(\bm{w}^{\mathsf T}\bm{x}_i+b)\ge 1.
  \label{eq:svm-primal}
\end{equation}
两条间隔边界之间的距离为 $2/\norm{\bm{w}}_2$，
因此最小化 $\norm{\bm{w}}_2$ 等价于最大化几何间隔\cite{cortes}.
拉格朗日对偶问题为
\begin{equation}
  \max_{\bm{\alpha}}
  \sum_{i=1}^{n}\alpha_i
  -\frac{1}{2}\sum_{i=1}^{n}\sum_{j=1}^{n}
  \alpha_i\alpha_j y_i y_j
  \bm{x}_i^{\mathsf T}\bm{x}_j,
\end{equation}
约束为 $\alpha_i\ge0$ 与 $\sum_i\alpha_i y_i=0$。
满足 $\alpha_i>0$ 的样本是支持向量，它们决定分类边界。

\subsubsection{软间隔、Hinge 损失与核技巧}

线性不可分时引入松弛变量：
\begin{equation}
  \min_{\bm{w},b,\bm{\xi}}
  \frac{1}{2}\norm{\bm{w}}_2^2
  +C\sum_{i=1}^{n}\xi_i,
  \quad
  y_i(\bm{w}^{\mathsf T}\bm{x}_i+b)\ge1-\xi_i,\ \xi_i\ge0.
\end{equation}
它等价于正则化 Hinge 损失
\begin{equation}
  \frac{1}{2}\norm{\bm{w}}_2^2+
  C\sum_{i=1}^{n}
  \max\{0,1-y_i f(\bm{x}_i)\}.
\end{equation}
参数 $C$ 越大，对间隔违例的惩罚越强。核技巧用
\begin{equation}
  K(\bm{x}_i,\bm{x}_j)
  =\phi(\bm{x}_i)^{\mathsf T}\phi(\bm{x}_j)
\end{equation}
替代显式高维映射。常见 RBF 核为
\begin{equation}
  K(\bm{x},\bm{x}')
  =\exp(-\gamma\norm{\bm{x}-\bm{x}'}_2^2).
\end{equation}
有效核矩阵应为对称半正定。SVM 适合中小规模、高维数据，
但核矩阵在大样本下计算和存储成本较高，$C$、$\gamma$ 也需要谨慎调节。

\subsection{随机森林}

\subsubsection{Bagging 与双重随机化}

随机森林以决策树为基学习器，同时进行样本随机化和特征随机化\cite{breiman}：
\begin{enumerate}
  \item 对原训练集进行有放回的 Bootstrap 抽样；
  \item 在每个树节点仅从随机特征子集中寻找最佳切分；
  \item 分类任务对各树投票，回归任务对各树输出取平均。
\end{enumerate}
若单棵树预测方差为 $\sigma^2$，树间平均相关系数为 $\rho$，
$T$ 棵树平均预测的方差近似为
\begin{equation}
  \mathrm{Var}(\bar f)
  =\rho\sigma^2+\frac{1-\rho}{T}\sigma^2.
\end{equation}
增加树数主要消除第二项，而随机特征子集通过降低 $\rho$ 进一步减少方差。

\subsubsection{袋外评估与局限}

单棵树的 Bootstrap 样本没有包含某个给定样本的概率约为
\begin{equation}
  \left(1-\frac{1}{n}\right)^n
  \longrightarrow e^{-1}\approx0.368.
\end{equation}
因此可用袋外样本（Out-of-Bag, OOB）直接估计泛化误差。
随机森林对非线性和特征交互有较强表达能力，对尺度变换通常不敏感；
但模型体积较大，整体解释性弱于单棵树。基于不纯度的特征重要性可能偏向
高基数或可切分点较多的特征，必要时应使用置换重要性进行复核。

\subsection{梯度提升与 AdaBoost}

\subsubsection{函数空间中的梯度下降}

梯度提升按顺序添加弱学习器，在函数空间中最小化经验风险\cite{friedman}：
\begin{equation}
  \min_F\sum_{i=1}^{n}\ell(y_i,F(\bm{x}_i)).
\end{equation}
第 $m$ 轮计算伪残差
\begin{equation}
  r_{im}
  =-\left.
  \frac{\partial\ell(y_i,F(\bm{x}_i))}
  {\partial F(\bm{x}_i)}
  \right|_{F=F_{m-1}},
\end{equation}
训练弱学习器 $h_m$ 拟合 $\{r_{im}\}$，再进行线性搜索：
\begin{align}
  \gamma_m
  &=\argmin_\gamma
  \sum_{i=1}^{n}
  \ell\!\left(y_i,F_{m-1}(\bm{x}_i)+
  \gamma h_m(\bm{x}_i)\right),\\
  F_m(\bm{x})
  &=F_{m-1}(\bm{x})+\eta\gamma_m h_m(\bm{x}).
\end{align}
平方损失下，伪残差就是 $y_i-F_{m-1}(\bm{x}_i)$；
一般损失下则是损失对当前预测的负梯度。

\subsubsection{正则化与 AdaBoost 联系}

学习率 $\eta$、弱学习器复杂度、迭代轮数、样本子采样和早停共同控制模型容量。
AdaBoost 使用指数损失 $\ell(y,F)=\exp[-yF]$，
其弱分类器权重为
\begin{equation}
  \alpha_m=\frac{1}{2}
  \log\frac{1-\varepsilon_m}{\varepsilon_m},
\end{equation}
样本权重按
\begin{equation}
  w_i^{(m+1)}
  \propto w_i^{(m)}
  \exp[-\alpha_m y_i h_m(\bm{x}_i)]
\end{equation}
更新，使被错误分类的样本在下一轮更受关注。
Gradient Boosting 的损失选择更灵活；XGBoost 等方法在此基础上进一步加入
二阶近似、显式树复杂度正则化和工程优化。

\begin{table}[H]
\centering
\caption{Bagging 与 Boosting 的主要区别}
\begin{tabularx}{\textwidth}{L{2.6cm}Y Y}
\toprule
\textbf{维度} & \textbf{Bagging / 随机森林} & \textbf{Boosting / 梯度提升}\\
\midrule
训练关系 & 多个学习器可并行训练 & 后一学习器依赖前一模型\\
主要作用 & 降低方差、提高稳定性 & 逐步降低偏差与目标损失\\
组合方式 & 平均或投票 & 加权求和\\
噪声敏感性 & 通常较稳健 & 可能持续拟合噪声，需正则化\\
\bottomrule
\end{tabularx}
\end{table}

\section{无监督学习模型}

\subsection{K-Means 聚类}

\subsubsection{目标函数与交替优化}

K-Means 将样本划分为 $K$ 个簇，最小化簇内平方和\cite{macqueen}：
\begin{equation}
  J(\{\mathcal{C}_k\},\{\bm{\mu}_k\})
  =\sum_{k=1}^{K}\sum_{\bm{x}_i\in\mathcal{C}_k}
  \norm{\bm{x}_i-\bm{\mu}_k}_2^2.
\end{equation}
Lloyd 算法交替执行两步：
\begin{align}
  c_i
  &\leftarrow\argmin_k
  \norm{\bm{x}_i-\bm{\mu}_k}_2^2,\\
  \bm{\mu}_k
  &\leftarrow\frac{1}{\abs{\mathcal{C}_k}}
  \sum_{\bm{x}_i\in\mathcal{C}_k}\bm{x}_i.
\end{align}
在中心固定时，最近中心分配最优；在簇分配固定时，样本均值使平方误差最小。
因此每轮迭代都不会增大目标函数，算法会在有限种划分中停止，
但只保证收敛到局部最优或稳定划分，而不保证全局最优。

\subsubsection{初始化、尺度与选取 $K$}

K-Means++ 先随机选一个中心，之后以样本到最近已有中心的平方距离为权重选择新中心，
可以减少劣质初始化。实践中还应多次随机运行并选择目标函数最小的结果。
由于欧氏距离受量纲影响，通常需要仅依据训练数据进行 Z-score 标准化：
\begin{equation}
  z_{ij}=\frac{x_{ij}-\mu_j}{\sigma_j}.
\end{equation}
$K$ 可结合手肘法、轮廓系数和领域知识确定。
K-Means 偏好大小相近的近球形簇，对异常值、非凸形状和初始中心较敏感。

\subsection{层次聚类}

层次聚类产生从细到粗的树状结构。凝聚式算法从每个样本单独成簇开始，
反复合并最近的两个簇；分裂式算法从一个总簇开始递归拆分。
常见簇间距离为
\begin{align}
  d_{\mathrm{single}}(A,B)
  &=\min_{\bm{x}\in A,\bm{y}\in B}d(\bm{x},\bm{y}),\\
  d_{\mathrm{complete}}(A,B)
  &=\max_{\bm{x}\in A,\bm{y}\in B}d(\bm{x},\bm{y}),\\
  d_{\mathrm{average}}(A,B)
  &=\frac{1}{\abs A\abs B}
  \sum_{\bm{x}\in A}\sum_{\bm{y}\in B}
  d(\bm{x},\bm{y}).
\end{align}
Ward 方法选择使类内平方和增量最小的合并：
\begin{equation}
  \Delta(A,B)
  =\frac{\abs A\abs B}{\abs A+\abs B}
  \norm{\bm{\mu}_A-\bm{\mu}_B}_2^2.
\end{equation}
单链接能识别链状结构但容易出现 chaining 效应；
完全链接产生较紧凑的簇；平均链接较折中；Ward 链接通常偏好近球形簇。
树状图可以在不同高度切割得到不同簇数，但早期错误合并通常无法撤销。
层次聚类至少需要 $O(n^2)$ 级别的距离存储，因而不适合超大样本。

\subsection{主成分分析}

\subsubsection{最大方差推导}

设中心化数据矩阵为
\begin{equation}
  \bm{X}_c
  =\bm{X}-\bm{1}\bm{\mu}^{\mathsf T},
  \qquad
  \bm{S}=\frac{1}{n-1}\bm{X}_c^{\mathsf T}\bm{X}_c.
\end{equation}
第一主成分方向通过最大化投影方差得到\cite{jolliffe}：
\begin{equation}
  \max_{\bm{w}}\ \bm{w}^{\mathsf T}\bm{S}\bm{w},
  \qquad
  \text{s.t.}\quad\bm{w}^{\mathsf T}\bm{w}=1.
\end{equation}
构造拉格朗日函数
\begin{equation}
  \mathcal{J}(\bm{w},\lambda)
  =\bm{w}^{\mathsf T}\bm{S}\bm{w}
  -\lambda(\bm{w}^{\mathsf T}\bm{w}-1),
\end{equation}
求导并令其为零：
\begin{equation}
  \bm{S}\bm{w}=\lambda\bm{w}.
\end{equation}
因此最优方向是最大特征值对应的特征向量。
后续主成分在与已有方向正交的约束下依次取次大特征值对应的特征向量。

\subsubsection{投影、解释方差与重建}

令特征值满足
$\lambda_1\ge\cdots\ge\lambda_d\ge0$，
$\bm{W}_k=[\bm{w}_1,\ldots,\bm{w}_k]$，则低维表示为
\begin{equation}
  \bm{Z}=\bm{X}_c\bm{W}_k,
\end{equation}
重建为
\begin{equation}
  \hat{\bm{X}}
  =\bm{Z}\bm{W}_k^{\mathsf T}
  +\bm{1}\bm{\mu}^{\mathsf T}.
\end{equation}
累计解释方差比为
\begin{equation}
  \mathrm{CVR}(k)
  =\frac{\sum_{j=1}^{k}\lambda_j}
  {\sum_{j=1}^{d}\lambda_j}.
\end{equation}
按每个样本归一化的最优重建误差为
$\sum_{j=k+1}^{d}\lambda_j$。
数值计算通常直接对 $\bm{X}_c$ 做 SVD，避免显式形成协方差矩阵。
PCA 是线性、无监督的方差保留方法；若高方差方向不是任务相关方向，
降维可能丢失重要信号。

\section{神经网络与表示学习模型}

\subsection{多层感知机}

多层感知机（MLP）由全连接层与非线性激活组成，其前向传播见式
\eqref{eq:layer}。若所有层都没有非线性激活，多层线性变换仍可合并成一个线性变换，
因此激活函数是表达复杂决策边界的关键。

\begin{table}[H]
\centering
\caption{常用激活函数及性质}
\begin{tabularx}{\textwidth}{L{2.2cm}L{4.7cm}Y}
\toprule
\textbf{函数} & \textbf{定义} & \textbf{主要特点}\\
\midrule
Sigmoid & $\sigma(z)=1/(1+e^{-z})$ & 输出在 $(0,1)$，饱和区易梯度消失\\
Tanh & $\tanh(z)=(e^z-e^{-z})/(e^z+e^{-z})$ & 零中心，但仍存在饱和问题\\
ReLU & $\max(0,z)$ & 计算简单，正半轴梯度稳定；可能出现死亡神经元\\
Softmax & $e^{z_k}/\sum_j e^{z_j}$ & 将多类 logits 转换为概率分布\\
\bottomrule
\end{tabularx}
\end{table}

为使信号方差在层间保持稳定，Xavier 初始化适合近似对称激活：
\begin{equation}
  \mathrm{Var}(W_{ij})
  \approx\frac{2}{n_{\mathrm{in}}+n_{\mathrm{out}}},
\end{equation}
ReLU 网络常使用 He 初始化：
\begin{equation}
  \mathrm{Var}(W_{ij})
  \approx\frac{2}{n_{\mathrm{in}}}.
\end{equation}
倒置 Dropout 在训练时采样
$m_j\sim\mathrm{Bernoulli}(q)$，并令
\begin{equation}
  \tilde a_j=\frac{m_j}{q}a_j,
\end{equation}
从而保持激活期望不变；推理时无需再缩放。
MLP 是理解前向传播、反向传播、初始化和正则化的基础模型，
但对图像等结构化输入使用全连接会产生大量参数，缺少空间归纳偏置。

\subsection{卷积神经网络}

卷积神经网络利用局部连接与权值共享处理网格状数据\cite{lecun}。
严格来说，深度学习框架中的“卷积”通常实现互相关。对二维输入，
\begin{equation}
  Y_{i,j,c_{\mathrm{out}}}
  =b_{c_{\mathrm{out}}}
  +\sum_{u=0}^{K_h-1}\sum_{v=0}^{K_w-1}
  \sum_{c=1}^{C_{\mathrm{in}}}
  W_{u,v,c,c_{\mathrm{out}}}
  X_{iS_h+u-P_h,\ jS_w+v-P_w,\ c}.
\end{equation}
输出空间尺寸为
\begin{align}
  H_{\mathrm{out}}
  &=\left\lfloor
  \frac{H_{\mathrm{in}}+2P_h-K_h}{S_h}
  \right\rfloor+1,\\
  W_{\mathrm{out}}
  &=\left\lfloor
  \frac{W_{\mathrm{in}}+2P_w-K_w}{S_w}
  \right\rfloor+1.
\end{align}
卷积层参数量为
\begin{equation}
  (K_hK_wC_{\mathrm{in}}+1)C_{\mathrm{out}},
\end{equation}
与输入图像的空间尺寸无关。池化或带步幅卷积用于降低空间分辨率，
扩大有效感受野。局部连接使网络优先学习边缘和纹理，
权值共享使相同模式可以在不同位置被检测。

从 LeNet-5 到 AlexNet、VGG、ResNet，CNN 的发展主要围绕网络深度、
优化稳定性和计算效率展开。CNN 对局部空间结构有强归纳偏置，
但远距离依赖需要通过堆叠多层或引入注意力机制建模。

\subsection{自注意力与 Transformer 基础}

知识库已把自注意力列入 Transformer 学习路径，但对应条目尚未展开\cite{aikb}。
给定输入矩阵 $\bm{X}\in\R^{n\times d}$，线性映射得到
\begin{equation}
  \bm{Q}=\bm{X}\bm{W}_Q,\qquad
  \bm{K}=\bm{X}\bm{W}_K,\qquad
  \bm{V}=\bm{X}\bm{W}_V.
\end{equation}
缩放点积注意力为\cite{vaswani}
\begin{equation}
  \mathrm{Attention}(\bm{Q},\bm{K},\bm{V})
  =\mathrm{softmax}\left(
  \frac{\bm{Q}\bm{K}^{\mathsf T}}{\sqrt{d_k}}
  +\bm{M}
  \right)\bm{V},
  \label{eq:attention}
\end{equation}
其中 $\bm{M}$ 是可选掩码。除以 $\sqrt{d_k}$ 是为了避免维度较大时点积方差过大，
使 Softmax 过度饱和。多头注意力在不同子空间独立计算：
\begin{align}
  \mathrm{head}_h
  &=\mathrm{Attention}
  (\bm{X}\bm{W}_Q^{(h)},
   \bm{X}\bm{W}_K^{(h)},
   \bm{X}\bm{W}_V^{(h)}),\\
  \mathrm{MHA}(\bm{X})
  &=\mathrm{Concat}(\mathrm{head}_1,\ldots,\mathrm{head}_H)
  \bm{W}_O.
\end{align}
自注意力本身不包含顺序信息，需要位置编码。
它可以直接建立任意位置之间的依赖，但标准实现的注意力矩阵大小为 $n\times n$，
时间和空间复杂度对序列长度均呈二次增长。

\section{模型比较与使用原则}

\small
\begin{longtable}{L{2.1cm}L{2.45cm}L{2.75cm}L{2.75cm}L{3.1cm}}
\caption{基本模型的统一比较}\label{tab:model-compare}\\
\toprule
\textbf{模型} & \textbf{核心假设或偏置} & \textbf{优化或求解} &
\textbf{优势} & \textbf{主要局限}\\
\midrule
\endfirsthead
\toprule
\textbf{模型} & \textbf{核心假设或偏置} & \textbf{优化或求解} &
\textbf{优势} & \textbf{主要局限}\\
\midrule
\endhead
\bottomrule
\endfoot
线性回归 & 条件均值近似线性 & OLS、梯度下降 & 简单、可解释、训练快 & 异常值和非线性敏感\\
逻辑回归 & 对数几率线性 & 最大似然、凸优化 & 概率输出、强 Baseline & 原空间决策边界线性\\
决策树 & 递归轴对齐划分 & 贪心划分与剪枝 & 可解释、能建模交互 & 高方差、对数据扰动敏感\\
朴素贝叶斯 & 类内特征条件独立 & 频率或分布参数估计 & 快、小样本友好 & 独立假设较强\\
SVM & 大间隔带来较好泛化 & 凸二次规划、核技巧 & 高维中小样本有效 & 大样本核计算昂贵\\
随机森林 & 多棵去相关树平均 & Bagging 与随机特征 & 稳定、鲁棒、少预处理 & 模型较大、整体解释性弱\\
梯度提升 & 加法模型逐步降低损失 & 函数梯度下降 & 精度高、损失灵活 & 串行训练、需防止拟合噪声\\
K-Means & 欧氏空间近球形簇 & 交替最小化 & 简单高效 & 需给定 $K$，局部最优\\
层次聚类 & 距离与链接反映层次 & 贪心合并或拆分 & 产生完整树状结构 & 计算和存储成本较高\\
PCA & 线性高方差方向含主要信息 & 特征分解或 SVD & 降维、去相关、可视化 & 无监督且只能线性投影\\
MLP & 组合非线性可逼近复杂映射 & 反向传播与梯度优化 & 通用表达能力强 & 参数多、结构先验弱\\
CNN & 局部性与平移等变性 & 反向传播 & 图像等网格数据高效 & 长程关系建模间接\\
自注意力 & 内容相关的全局交互 & 反向传播 & 直接建模长距离依赖 & 标准复杂度为 $O(n^2)$\\
\end{longtable}
\normalsize

选择模型时不应只根据“先进程度”。更合理的顺序是：
\begin{enumerate}
  \item 明确任务是回归、分类、聚类、降维还是序列决策；
  \item 分析样本量、特征类型、标注质量、计算预算与解释需求；
  \item 从线性或树模型建立可复现 Baseline；
  \item 依据误差分析判断是否需要非线性、集成或深度表示；
  \item 在同一数据划分和指标下公平比较，并报告训练成本与局限。
\end{enumerate}

\section{总结与后续学习计划}

通过本次整理，我将机器学习入门知识统一为六个相互关联的问题：
\begin{enumerate}
  \item 数据和任务如何定义；
  \item 模型引入了什么结构假设；
  \item 损失函数如何把目标转化为可优化量；
  \item 前向传播、反向传播和优化器分别承担什么作用；
  \item 什么条件可以视为收敛，收敛是否能够泛化；
  \item 应使用哪些指标和实验设计支持最终结论。
\end{enumerate}
线性回归、逻辑回归、树模型、概率模型、间隔模型、集成模型、
聚类、降维和神经网络虽然形式不同，但都能放入
“表示--目标--求解--评价”的框架中理解。

下一步我将把理论学习与代码验证结合：首先使用统一数据集实现线性回归、
逻辑回归和决策树 Baseline；随后比较随机森林与梯度提升；
再用 K-Means 和 PCA 观察无监督结构；最后用 PyTorch 实现 MLP 和 CNN 的
前向与反向传播，并记录不同学习率、初始化方式和正则化设置对收敛曲线的影响。
每次实验都保存数据划分、随机种子、配置、日志和评价指标，
使数学理解能够通过可复现实验得到检验。

%==================== 参考文献 ====================
\clearpage
\small
\raggedright
\begin{thebibliography}{99}

\bibitem{aikb}
CCYSAM.
AI-Knowledge-Base: Machine Learning and Deep Learning[EB/OL].
(2026-07-06)[2026-07-29].
\url{https://git.ccysam.xyz/ccysam/AI-Knowledge-Base}.

\bibitem{mitchell}
MITCHELL T M.
Machine Learning[M].
New York: McGraw-Hill, 1997.

\bibitem{goodfellow}
GOODFELLOW I, BENGIO Y, COURVILLE A.
Deep Learning[M].
Cambridge, MA: MIT Press, 2016.

\bibitem{sutton}
SUTTON R S, BARTO A G.
Reinforcement Learning: An Introduction[M].
2nd ed.
Cambridge, MA: MIT Press, 2018.

\bibitem{bishop}
BISHOP C M.
Pattern Recognition and Machine Learning[M].
New York: Springer, 2006.

\bibitem{hastie}
HASTIE T, TIBSHIRANI R, FRIEDMAN J.
The Elements of Statistical Learning: Data Mining, Inference, and Prediction[M].
2nd ed.
New York: Springer, 2009.

\bibitem{rumelhart}
RUMELHART D E, HINTON G E, WILLIAMS R J.
Learning representations by back-propagating errors[J].
Nature, 1986, 323: 533--536.
DOI:10.1038/323533a0.

\bibitem{kingma}
KINGMA D P, BA J.
Adam: A method for stochastic optimization[C]//
Proceedings of the 3rd International Conference on Learning Representations.
San Diego: ICLR, 2015.

\bibitem{cortes}
CORTES C, VAPNIK V.
Support-vector networks[J].
Machine Learning, 1995, 20: 273--297.
DOI:10.1007/BF00994018.

\bibitem{breiman}
BREIMAN L.
Random forests[J].
Machine Learning, 2001, 45: 5--32.
DOI:10.1023/A:1010933404324.

\bibitem{friedman}
FRIEDMAN J H.
Greedy function approximation: A gradient boosting machine[J].
The Annals of Statistics, 2001, 29(5): 1189--1232.
DOI:10.1214/aos/1013203451.

\bibitem{macqueen}
MACQUEEN J.
Some methods for classification and analysis of multivariate observations[C]//
Proceedings of the Fifth Berkeley Symposium on Mathematical Statistics and Probability.
Berkeley: University of California Press, 1967: 281--297.

\bibitem{jolliffe}
JOLLIFFE I T, CADIMA J.
Principal component analysis: A review and recent developments[J].
Philosophical Transactions of the Royal Society A, 2016, 374(2065): 20150202.
DOI:10.1098/rsta.2015.0202.

\bibitem{lecun}
LECUN Y, BOTTOU L, BENGIO Y, et al.
Gradient-based learning applied to document recognition[J].
Proceedings of the IEEE, 1998, 86(11): 2278--2324.
DOI:10.1109/5.726791.

\bibitem{vaswani}
VASWANI A, SHAZEER N, PARMAR N, et al.
Attention is all you need[C]//
Advances in Neural Information Processing Systems 30.
Long Beach: Curran Associates, 2017: 5998--6008.

\end{thebibliography}
\normalsize

\end{document}
